\section{Introduction}
The idea of studying the Kepler problem, or the problem of $n$ bodies, in a non-Euclidean space has a long history, but in almost all cases such
generalisations were performed in spaces of~con\-stant curvature. For a very detailed and critical recent overview
of this subject we refer to the nice article \cite{borisov:16::}.

In \cite{1} the authors considered the question of generalization of
the Kepler problem and its resulting mechanics, based on first principles. The idea was to
formulate an analog of the classical problem in spaces which are
homogeneous, isotropic and admit dilations. The last requirement is
very restrictive because among homogeneous Riemannian manifolds only
Euclidean ones admit dilations. The simplest non-Euclidean metric
space satisfying all the required properties is the Heisenberg group
of special upper-triangular matrices
\begin{equation*}
 \begin{pmatrix}
 1 & x_1 & x_3 \\ 0 & 1 & x_2 \\ 0 & 0 & 1
 \end{pmatrix}\!,
 \qquad x_i\in\mathbb{R}.
\end{equation*}
An isomorphic representation is obtained by taking new coordinates
\begin{equation*}
 x_1 = x, \qquad x_2 = y, \qquad x_3 = z +\frac12 x_1 x_2,
\end{equation*}
in which the group action is
\begin{equation}
 (x_1,y_1,z_1)\cdot(x_2,y_2,z_2)=\bigg(x_1+x_2,y_1+y_2,z_1+z_2+\frac12(x_1 y_2-x_2 y_1)\bigg).
 \label{groupop}
\end{equation}

This manifold carries a maximally non-integrable distribution spanned
by the vector fields
\begin{equation*}
 X_1 = \partial_x - \frac12 y\partial_z,\qquad
 X_2 = \partial_y + \frac12 x\partial_z,
\end{equation*}
and which define the sub-Riemannian structure. Furthermore, there
exists a sub-elliptic Laplacian $\Delta:=X_1^2+X_2^2$ for which the
fundamental solution to the Poisson equation $\Delta U=\varrho$, with density $\varrho$, is known~\cite{Folland}. This makes the analogy complete because one can
construct the kinetic energy with the sub-Riemannian metric, and take
the potential as the point-source solution $U$;
the Hamiltonian of such Kepler--Heisenberg system is then
\begin{equation}
 \label{KH}
 H = \frac12 \bigg(p_x - \frac12 y p_z\bigg)^2 + \frac12\bigg(p_y+\frac12 x p_z\bigg)^2
 - \frac{\kappa}{\rho}, \qquad
 \rho:=\sqrt{(x^2+y^2)^2+{16}z^2},
\end{equation}
where $\kappa$ is a non-zero real parameter. The properties of this
system have been analysed in \cite{dots, 1,2}, where it was shown that it
has an invariant submanifold given by $z=0$, $p_z=0$,
$p_{\theta} = x p_y -y p_x=0$ on which all trajectories are straight
lines. More importantly, it was also demonstrated that all periodic
solutions must lie on the zero-energy level and that the whole system
is Liouville integrable there. This happens because of the quantity
$J=x p_x+ y p_y +2 z p_z$, for which $\dot{J} = 2H$, and on the level
$H=0$, it is the third integral of motion next to $H$ and
$p_{\theta}$, with which it commutes.
\begin{Remark}\label{rem1}
 The form of potential in~\eqref{KH} is the same as in~\cite{dots} which is the correction of that~in~\cite{1}.
\end{Remark}
We complete the above findings by proving that for all the other
values of energy, the system is not integrable. In addition, the generalisation of the above system to two bodies is straightforward, and we prove its non-integrability as well.
